% =============================================================================
% base.tex —— 基础层：宏包、中文字体、表格、文献、正文用自定义命令
% -----------------------------------------------------------------------------
% 加载顺序见 main.tex；本文件在 Beamer 类之后、其余公共配置之前加载。
% 本文件只依赖「已经载入 beamer」这一前提，不依赖配色与版式设置。
% =============================================================================

% -----------------------------------------------------------------------------
% 1. 中文与字体
% -----------------------------------------------------------------------------
% Tectonic / XeLaTeX 编译；不让 fontspec 接管数学字体，避免与 amssymb 冲突。
% ctex 在 fontset=none 下只做中文排版机制，字体由下面的 fontspec 命令显式指定。
\PassOptionsToPackage{no-math}{fontspec}
\usepackage[UTF8,fontset=none]{ctex}
\defaultfontfeatures{Ligatures=TeX}

% 字体策略：优先使用随项目携带的思源宋体/黑体（fonts/ 目录，无需系统安装）；
% 如果这些 OTF 被删掉，则回退到 TeX 发行版自带的 Fandol。
% 四个文件必须全部存在，避免字体的 Regular / Bold 指向缺失文件。
% 中文没有西文意义上的 italic：斜体统一映射到 Fandol 楷体，
% 粗斜体保持黑体/宋体粗体，避免出现伪斜体。
\newif\ifInkBundledFonts
\InkBundledFontstrue
\newcommand{\InkCheckFont}[1]{\IfFileExists{fonts/#1}{}{\InkBundledFontsfalse}}
\forcsvlist{\InkCheckFont}{SourceHanSerifCN-Regular.otf,SourceHanSerifCN-Bold.otf,SourceHanSansCN-Regular.otf,SourceHanSansCN-Bold.otf}
\ifInkBundledFonts
  % 完整思源字体：衬线、无衬线与中文等宽族
  \setCJKmainfont{SourceHanSerifCN-Regular.otf}[Path=fonts/,
    BoldFont=SourceHanSerifCN-Bold.otf,
    ItalicFont=FandolKai-Regular.otf,ItalicFeatures={Path=},
    BoldItalicFont=SourceHanSerifCN-Bold.otf]%
  \setCJKsansfont{SourceHanSansCN-Regular.otf}[Path=fonts/,
    BoldFont=SourceHanSansCN-Bold.otf,
    ItalicFont=FandolKai-Regular.otf,ItalicFeatures={Path=},
    BoldItalicFont=SourceHanSansCN-Bold.otf]%
  % 等宽场景（代码、路径）用思源黑体，保证中文可读
  \setCJKmonofont{SourceHanSansCN-Regular.otf}[Path=fonts/,
    BoldFont=SourceHanSansCN-Bold.otf,
    ItalicFont=FandolKai-Regular.otf,ItalicFeatures={Path=},
    BoldItalicFont=SourceHanSansCN-Bold.otf]%
\else
  % 思源字体不完整：整套回退到 Fandol
  \setCJKmainfont{FandolSong-Regular.otf}[BoldFont=FandolSong-Bold.otf,
    ItalicFont=FandolKai-Regular.otf,BoldItalicFont=FandolSong-Bold.otf]%
  \setCJKsansfont{FandolHei-Regular.otf}[BoldFont=FandolHei-Bold.otf,
    ItalicFont=FandolKai-Regular.otf,BoldItalicFont=FandolHei-Bold.otf]%
  \setCJKmonofont{FandolFang-Regular.otf}[BoldFont=FandolHei-Bold.otf,
    ItalicFont=FandolKai-Regular.otf,BoldItalicFont=FandolHei-Bold.otf]%
\fi
% 西文保留 Latin Modern（beamer 默认），无需额外设置。

% -----------------------------------------------------------------------------
% 2. 表格
% -----------------------------------------------------------------------------
% booktabs 提供 \toprule/\midrule/\bottomrule 三线表；
% tabularx 提供自适应列宽，下面定义 Y 列 = 自动宽度 + 左对齐（避免两端对齐的稀疏感）。
\usepackage{booktabs,tabularx}
\newcolumntype{Y}{>{\raggedright\arraybackslash}X}
% 表头按单元格显式标记；普通表格、自动宽度表格与合并单元格均可复用。
\newcommand{\tablehead}[1]{\textbf{#1}}
% 表格统一使用衬线体：中文为思源宋体、西文为 Latin Modern Roman。
% 用环境分组包住，保证表外的正文与来源说明仍用原字体。
\newcommand{\InkTableStyle}{\rmfamily\renewcommand{\arraystretch}{1.2}}
\AtBeginEnvironment{tabular}{\InkTableStyle}
\AtBeginEnvironment{tabularx}{\InkTableStyle}

% -----------------------------------------------------------------------------
% 3. 图片与链接
% -----------------------------------------------------------------------------
% 图片统一放在 figures/，正文里只写文件名即可
\graphicspath{{figures/}}
% 网址与正文用同一套字形，避免 URL 突然换字体
\urlstyle{same}

% -----------------------------------------------------------------------------
% 4. 引用类文字的专用字体（Times 风格）
% -----------------------------------------------------------------------------
% 参考文献、引文块与脚注使用 Times 风格衬线体，
%       而正文西文仍保留 Latin Modern。
% 说明：TeX 发行版不自带 Times New Roman（商业字体），这里使用其度量兼容的
%       TeX Gyre Termes（Nimbus Roman No9 的再发布版），字形与 Times New Roman
%       基本一致，且随 Tectonic / TeX Live 一起分发，无需额外安装。
% 使用文件名加载，避免依赖系统字体索引。
% 若某套发行版只提供 X 变体，则回退到 TeXGyreTermesX-*.otf，保证仍可编译。
\IfFileExists{texgyretermes-regular.otf}{%
  \newfontfamily\TimesCite{texgyretermes-regular.otf}[%
    BoldFont=texgyretermes-bold.otf,
    ItalicFont=texgyretermes-italic.otf,
    BoldItalicFont=texgyretermes-bolditalic.otf]%
}{%
  \newfontfamily\TimesCite{TeXGyreTermesX-Regular.otf}[%
    BoldFont=TeXGyreTermesX-Bold.otf,
    ItalicFont=TeXGyreTermesX-Italic.otf,
    BoldItalicFont=TeXGyreTermesX-BoldItalic.otf]%
}

% 引用块环境：沿用 LaTeX 标准名 quote / quotation，但换成 Times 风格字体。
% quote 适合较短引文（整段缩进、无首行缩进）；
% quotation 适合多段落长引文（整段缩进且首行再缩进）。
% 保存原始环境命令，包裹字体切换并保持原有缩进和分段规则。
\let\origquote\quote
\let\endorigquote\endquote
\let\origquotation\quotation
\let\endorigquotation\endquotation
\renewenvironment{quote}
  {\origquote\TimesCite\small}
  {\endorigquote}
\renewenvironment{quotation}
  {\origquotation\TimesCite\small}
  {\endorigquotation}

% -----------------------------------------------------------------------------
% 5. 参考文献
% -----------------------------------------------------------------------------
% 后端固定为 BibTeX：Tectonic 内置 BibTeX，但**不提供 biber**，
% 所以这里不能改成 backend=biber（会以 "program not found" 失败）。
% 若将来需要 biber 的完整多语种/中文作者支持，必须改用带 biber 的
% XeLaTeX 工具链（见 README「参考文献」一节）。
% 作者年份制、按姓名-年份-标题排序，正文引用最多 2 位作者，文献表最多 6 位。
\usepackage[backend=bibtex,style=authoryear,sorting=nyt,
  maxcitenames=2,maxbibnames=6,giveninits=true,uniquename=init,
  doi=true,url=true,isbn=false]{biblatex}
\addbibresource{references/references.bib}

% 参考文献表字体：必须通过 biblatex 的 \bibfont 钩子设置。
% 在 \bibfont 中设置，避免条目排版时重置回正文字体。
\renewcommand{\bibfont}{\TimesCite}

% csquotes：提供语言自适应的引号命令 \enquote{...} 与长引文 \blockquote{...}，
% 免去手工写 ``...'' 这类 TeX 引号；autostyle 让它跟随文档语言选择引号形式。
\usepackage[autostyle]{csquotes}

% -----------------------------------------------------------------------------
% 6. 正文中使用的可复用命令
% -----------------------------------------------------------------------------
% 突出重点：使用主题结构色 + 粗体。
\newcommand{\keypoint}[1]{\structure{\textbf{#1}}}

% 独立段落注释：无标题栏的浅灰圆角盒，小一号灰色直立正文。
% 原生圆角盒左右各外扩 4bp，内缩后与当前栏边界对齐。
\newcommand{\slideNote}[1]{%
  \par\smallskip\begingroup
  \noindent\hspace*{4bp}%
  \begin{minipage}{\dimexpr\linewidth-8bp\relax}
    \begin{beamercolorbox}[wd=\linewidth,sep=0.65em,rounded=true,shadow=false]{annotation}
      \small\sffamily\mdseries\upshape\raggedright
      \setlength{\parskip}{0.25em}#1\par
    \end{beamercolorbox}
  \end{minipage}\par\endgroup\smallskip}

% 来源说明（图表出处、数据来源）：小字号、灰色、左对齐，无脚注编号。
% 内容通常是 \parencite 之类的引用，因此一并使用 Times 风格字体。
% 在调用处就地排版，建议写在 frame 末尾。
\newcommand{\sourceNote}[1]{%
  \par\smallskip
  {\raggedright\SourceNoteSize\TimesCite\color{AcademicGray}#1\par}}

% 插图：第一参数为高度上限（默认 0.58\textheight），宽度自动适应单栏或双栏，
% 并保持原始比例。例：\slideimage[0.45\textheight]{result.png}
\newcommand{\slideimage}[2][0.58\textheight]{%
  \includegraphics[width=\linewidth,height=#1,keepaspectratio]{#2}}
